\chapter{Discussion}
\label{ch:discussion}

\section{Hypothesis Evaluation}
\label{sec:disc_hypotheses}

\subsection*{H1: Non-Convex Shapes Have Lower Success Rates Than Convex Shapes}

H1 is confirmed strongly. The Circle and Square achieved aggregate success rates of 88.1\% and 83.9\% respectively, while the L-shape achieved 62.4\% and the C-shape 23.0\%. The gap between the best-performing and worst-performing shapes is 65.1 percentage points, and the ordering is monotonically related to geometric complexity. The convex-versus-non-convex divide is not a marginal effect: it is the single largest source of variance in the dataset, dwarfing the contributions of swarm size and shuffle randomness within any given shape condition. The C-shape result in particular — a success rate more than twice as low as the next worst shape — underscores that geometric concavity introduces a qualitatively different transport problem, not merely a harder instance of the same problem.

\subsection*{H2: An Optimal Intermediate Shuffle Randomness Level Exists for Non-Convex Shapes}

H2 is partially confirmed for the L-shape but is not confirmed for the C-shape. For the L-shape, the relationship between SR and success rate is monotonically decreasing: SR = 0.0 (fully structured orbital redistribution) achieves the highest success rate and SR = 1.0 the lowest, with intermediate values falling between these extremes. The optimal SR level is therefore at the structured extreme rather than at an intermediate value, which partially contradicts H2's prediction of an intermediate optimum. This suggests that, for the L-shape, the orbital redistribution mechanism is superior not merely to full randomness but to any mixture of the two strategies within the tested range.

For the C-shape, no clear relationship between SR and success rate was detected across the 0.0--0.75 range in the main experiment. The coarse five-level sweep may have been insufficient to resolve the narrow optimal band that the bi-stability literature predicts for shapes with deep concavities \citep{springer2022}. The follow-up experiment, which employs an 11-level SR sweep at finer resolution, is designed to address this limitation directly.

\subsection*{H3: Success Rate Plateaus or Declines at Large Swarm Sizes for Non-Convex Shapes}

H3 is confirmed. Both the L-shape and the C-shape exhibit declining performance at swarm size 35 relative to smaller sizes, while the convex shapes show monotonic improvement across the full size range. The C-shape presents the most pronounced effect: success rate peaks at size 25 and declines monotonically thereafter, with size 35 performing worse than size 25 at every SR level tested. The failure mode analysis corroborates this interpretation: at size 35, C-shape failures are almost exclusively of the stuck type, with 196 out of 198 failures classified as stuck failures, indicating that agents reliably locate and attach to the payload but the collective fails to complete transport. This pattern is consistent with the hypothesis that additional agents beyond the optimum disproportionately populate the interior concavity of the C-shape, where their applied forces are geometrically counterproductive. The finding has practical implications for swarm sizing in applications involving non-convex objects, where deploying larger swarms may actively worsen outcomes rather than improving them.

\subsection*{H4: Peak\_Attached Mediates the Relationship Between Swarm Size and Success}

H4 is supported directionally but requires formal mediation analysis for confirmation. Peak\_Attached increases with swarm size across all shape conditions, consistent with the proposed role of this variable as a mediator. For the convex shapes, higher Peak\_Attached values are associated with higher success rates, which is consistent with the straightforward interpretation that more simultaneously attached agents generate greater collective force and increase the probability of successful transport. For the C-shape at large swarm sizes, however, Peak\_Attached continues to rise while success rate falls. This divergence is consistent with the proposed mediation mechanism: attachments occurring within the concavity contribute to Peak\_Attached but reduce rather than increase the net effective force applied to the payload. The result suggests that Peak\_Attached is a useful mediating variable for convex shapes but must be interpreted with caution for concave shapes, where raw attachment counts conflate productive and counterproductive attachments. A full formal mediation analysis, including Baron-Kenny steps or a bootstrap indirect-effects test, is deferred pending the availability of the complete follow-up data, which records per-trial attachment process variables that will enable more refined analyses.

\section{The Geometric Deadlock Problem}
\label{sec:disc_deadlock}

The C-shape results provide clear evidence of geometry-induced failure that cannot be overcome by adding agents or by tuning the shuffling randomness parameter within the range tested. The concavity of the C-shape acts as a force sink: agents that attach to the inner surface of the shape apply inward-pointing forces directed roughly toward the centroid of the concavity rather than toward the home base. These inward forces partially cancel the outward forces applied by agents on the outer convex surface, reducing the net collective force available for transport. This interpretation is consistent with the formal analysis in \citet{arxiv2024concave}, who demonstrate that agents positioned within concavities are geometrically counterproductive for collective transport and that the probability of concavity population increases with swarm size, providing a mechanistic explanation for the non-monotonic size-success relationship observed here.

The finding that C-shape performance degrades with swarm size rather than simply plateauing is particularly notable. A plateau would suggest that additional agents are neutral at large swarm sizes — that the benefits of additional force are offset by the costs of concavity population, reaching an equilibrium. The observed decline instead indicates that the costs of concavity population scale super-linearly with swarm size within the range tested, making the problem actively worse at the larger sizes. Whether this trend would continue beyond size 35 or would eventually plateau at still larger swarm sizes is an open empirical question that the current design does not address.

\section{Shuffle Duration as the Primary Lever}
\label{sec:disc_duration}

Preliminary pilot results from the follow-up experiment suggest that shuffle duration, rather than shuffle randomness, is the variable that most meaningfully affects stuck-failure resolution for non-convex shapes. The pilot observed a success rate increase from approximately 74\% at a shuffle duration of 1 second to approximately 94\% at 8 seconds for the C-shape — a gain of 20 percentage points achieved by varying duration alone, compared to the negligible SR effect observed in the main experiment. This finding is consistent with \citet{wilson2014}, who identify detachment duration as the primary tunable parameter governing deadlock escape in collective transport systems, and with \citet{springer2022}, who show that longer repositioning periods specifically benefit shapes with deep deadlocks by allowing agents sufficient time to traverse to attachment positions on the geometrically productive outer surface.

The mechanistic interpretation is straightforward: agents that detach following a jam event and then re-attach after a short shuffle duration may not have moved far enough from their previous attachment position to materially alter the force geometry. In particular, for the C-shape, an agent that detaches from a position inside the concavity but shuffles for only one second may re-attach in a position that is still inside or near the concavity. Longer shuffle durations allow agents to travel further, increasing the probability that they re-attach to the outer productive surface. This interpretation implies a spatial traversal hypothesis: the minimum effective shuffle duration is determined by the time required for an agent to traverse from the deepest concavity position to the nearest outer surface position, which is a function of payload geometry and agent speed.

These findings are based on a two-replicate pilot and should be treated as directional rather than definitive. The full 50-replicate run will enable formal analysis of the duration-SR interaction and confirmation of the monotonic duration-success trend.

\section{Stochasticity as a Double-Edged Mechanism}
\label{sec:disc_stochasticity}

Across the main experiment, SR = 1.0 consistently produced the worst or near-worst outcomes for non-convex shapes. For the L-shape, SR = 1.0 achieved the lowest success rate at every swarm size. For the C-shape, SR = 1.0 produced the greatest performance degradation in the follow-up pilot, substantially below all lower SR levels. Full randomness during the shuffling phase destroys the orbital structure that allows agents to distribute evenly around the payload perimeter during repositioning. When agents perform a fully random Brownian walk, their re-attachment positions are effectively drawn from whatever distribution they happen to reach in the available time, which for the C-shape disproportionately re-populates the concavity simply because the concavity subtends a large angular fraction of the payload's boundary.

At intermediate SR levels — 0.1 through 0.75 — performance differences for the C-shape were small in the main experiment, consistent with the bi-stability hypothesis that there exists a narrow optimal randomness band but that the coarse five-level sweep was insufficient to resolve it \citep{springer2022}. The 11-level sweep employed in the follow-up is designed to characterise this region at higher resolution. For the L-shape, the monotonically decreasing SR-success relationship suggests that bi-stability, if present, would only manifest at SR levels below 0.0 if such a level were achievable — that is, the optimal randomness level is at or below the most structured end of the tested range. This result is consistent with the L-shape's simpler concavity: a single right-angle recess that does not produce the symmetric force-cancellation problem observed in the C-shape, and for which purely structured orbital redistribution is sufficient to resolve the majority of deadlocks.

\section{Limitations}
\label{sec:disc_limitations}

Several limitations of the present study should be acknowledged. First, the simulation employs a simplified two-dimensional rigid-body physics model. Real ant workers have flexible bodies, variable grip strength dependent on substrate and orientation, and the capacity to position themselves in three dimensions relative to a transported object. These factors influence the geometry of force application in ways that the flat 2D model does not capture. Results regarding the relationship between geometric concavity and collective transport failure may not transfer directly to physical experiments without appropriate validation.

Second, the jam detection threshold is fixed at a payload speed of 5.0 pixels per second averaged over a rolling 60-frame window. This threshold is not adaptive to payload mass, shape, or swarm size. A condition with a large swarm and a heavy payload may exhibit genuine productive transport at speeds below this threshold, causing unnecessary jam events to be triggered; conversely, a condition with a small swarm and a light payload may exhibit genuine stuck behaviour at speeds slightly above the threshold, causing jams to go undetected. An adaptive threshold that calibrates to condition-specific dynamics would improve the fidelity of the jam detection mechanism.

Third, the main experiment employed only five SR levels (0.0, 0.25, 0.5, 0.75, 1.0). For the C-shape, no clear SR effect was observed across the 0.0--0.75 range, but the coarse resolution means that a narrow optimal band — if one exists — would not have been detected. The follow-up experiment addresses this directly with an 11-level sweep.

Fourth, the follow-up pilot results are based on two replicates per condition. Two replicates are statistically underpowered for estimating condition means reliably, particularly in the presence of the high trial-to-trial variance characteristic of stochastic multi-agent simulations. The pilot findings should be treated as directional indicators pending the full 50-replicate run.

Fifth, the pheromone model uses zero diffusion across all layers, meaning pheromone values do not spread to neighbouring grid cells. This design choice preserves trail sharpness and directional information but sacrifices biological realism: real ant pheromones diffuse spatially, creating gradient fields that agents can sense from a distance without requiring direct contact with a previously visited cell. The no-diffusion model may underestimate the role of chemical communication in coordinating agents at distances greater than the agent vision radius, and findings regarding pheromone-mediated recruitment may not generalise to models with spatially spreading gradients.

\section{Future Work}
\label{sec:disc_future}

The most immediate priority is the completion of the follow-up experiment at full replicate count (13,200 trials) and the formal analysis of the shuffle duration by SR interaction. Specifically, the question of whether the duration-success relationship is mediated by SR level — that is, whether the benefit of longer durations is modulated by the degree of randomness during repositioning — requires the statistical power that the full run will provide.

Beyond the follow-up, several directions merit investigation. Adaptive jam detection that calibrates the trigger threshold to swarm size, payload shape, and current transport velocity would improve the precision of the intervention mechanism and could be incorporated into the existing simulation architecture without major restructuring. Extension to a three-dimensional simulation environment would assess whether the planar findings regarding concavity-induced force cancellation generalise when agents can position themselves in a third dimension relative to the payload surface. Finally, the explicit concavity-avoidance attachment policy proposed by \citet{arxiv2024concave} — in which agents are discouraged from attaching to identified concave surface regions — represents a targeted algorithmic intervention that could be implemented in the existing FSM framework and tested against the baseline shuffling mechanism for non-convex shapes.
